% !PRINTED 41
\weberpage{41}

Setzt man hierin $B_1, B_2 \ldots B_\mu$ für $\Theta$, so folgt, dass $p_1 m$
durch jede der Zahlen $b_1, b_2 \ldots b_\mu$ theilbar ist. Es ist also $p_1$
theilbar durch $p_1'$ und $m$ durch das kleinste gemeinschaftliche
Vielfache von $b_1', b_2' \ldots b_\mu'$. In diesem Schlusse können wir nun
durchweg $A$ mit $B$ vertauschen. Es muss also auch $p_1'$ durch $p_1$
theilbar sein, und daher
\weq{5}{p_1 = p_1',}
und ausserdem ergiebt sich, dass $m$ auch das kleinste gemeinschaftliche
 Vielfache von $b_1', b_2' \ldots b_\mu'$ ist.

Um daraus unseren Satz allgemein zu beweisen, nehmen wir an,
es sei für irgend ein $s$ bewiesen, dass
\weq{6}{p_1 = p_1', p_2 = p_2', \ldots p_{s-1} = p_{s-1}'}
sein müsse. Nach (4) ist für jedes Element $\Theta$:
\weq{7}{\Theta^{p_s m} = A_1^{\alpha_1 p_s m} A_2^{\alpha_2 p_s m} \ldots A_{s-1}^{\alpha_{s-1} p_s m},}
worin $m$ wie oben das kleinste gemeinschaftliche Vielfache von
$a_1', a_2' \ldots a_\nu'$ und zugleich von $b_1', b_2' \ldots b_\mu'$, also durch $p$ nicht
theilbar ist.

Wir bestimmen nun die Anzahl der in der Form (7) enthaltenen
 von einander verschiedenen Elemente.

Lassen wir $\alpha_1$ die Reihe der Zahlen $0, 1, 2 \ldots \dfrac{p_1}{p_s} - 1$
durchlaufen, so sind die Elemente
\weq{8}{1,\ A_1^{p_s m},\ A_1^{2 p_s m} \ldots A_1^{\left(\frac{p_1}{p_s}-1\right) p_s m}}
alle von einander verschieden, während $A_1^{\frac{p_1}{p_s} p_s m}$ wieder $= 1$
wird, so dass sich alle anderen Potenzen $A_1^{\alpha_1 p_s m}$ in der Reihe (8)
wiederfinden. Ebenso schliessen wir in Bezug auf die übrigen
Factoren von (7), woraus man die genaue Anzahl der von einander
 verschiedenen in $\Theta^{p_s m}$ enthaltenen Elemente
\weq{9}{z = \frac{p_1}{p_s} \frac{p_2}{p_s} \ldots \frac{p_{s-1}}{p_s}}
findet.

Von den beiden Zahlen $p_s$, $p_s'$ wird, wenn sie nicht gleich
sind, eine die grössere sein. Wir wollen also annehmen, es sei
\weq{10}{p_s' \geqq p_s.}

Nun drücken wir $\Theta$ durch die Basis $B$ aus, und setzen
\weq{11}{\Theta^{p_s m} = B_1^{\beta_1 p_s m} B_2^{\beta_2 p_s m} \ldots B_\mu^{\beta_\mu p_s m},}
%CONT>>
